\documentclass[varwidth=15cm,border=10pt]{standalone}
\usepackage{amsmath,xcolor,array}
\usepackage{fontspec,unicode-math}
\setmainfont{DejaVu Sans}[ItalicFont={DejaVu Sans},ItalicFeatures={FakeSlant=0.2},BoldItalicFont={DejaVu Sans Bold},BoldItalicFeatures={FakeSlant=0.2}]
\setmathfont{texgyredejavu-math.otf}
\setmathfont{DejaVu Sans}[range={up,bfup}]
\setmathfont{DejaVu Sans}[range={it,bfit},FakeSlant=0.2]
\setmathfont{texgyredejavu-math.otf}[range={"2202,"1D715,"03B5,"03F5,"2112,"2131}]
\pagecolor[HTML]{1E1E2E}\color[HTML]{F4F4F4}
\definecolor{accent}{HTML}{89B4FA}
\definecolor{warn}{HTML}{F38BA8}
\definecolor{good}{HTML}{A6E3A1}
\definecolor{ref}{HTML}{9399B2}
\newcommand{\heading}[1]{\par\medskip{\large\bfseries\color{accent}#1}\par\smallskip}
\newcommand{\src}[1]{{\color{ref}\small[#1]}}
\newcommand{\dolabel}{{\bfseries DO}\ \ }
\newcommand{\why}{{\bfseries\color{accent}WHY}\ \ }
\newcommand{\chk}{{\bfseries\color{good}CHECK}\ \ }
\newcommand{\rep}{{\bfseries\color{warn}WRITE}\ \ }
\setlength{\parindent}{0pt}
\raggedright
\begin{document}

{\Large\bfseries Solving A3 Question 1 yourself — the procedure}

\heading{Before you start}
L05 slide 29, on this assignment: {\itshape ``The two settings do not have the
same answer to `which method?', and saying which, and {\bfseries why}, is what is
being assessed.''} So the deliverable is an argument, and the numbers are its
evidence. Every step below carries what to do, why it is right, how to know it
worked, and what belongs in the report.

{\bfseries The one idea.} There is no formula for $p$, so you cannot
differentiate it \src{slides 3--4}. But $A(x)p(x)=q$ is an {\itshape identity},
true for every $x$, so differentiate {\itshape that} \src{slide 17}. And since
$A=\sum_e x_eA_e$ is linear, $\partial A/\partial x_e=A_e$ free of charge
\src{slide 11}. Everything else is bookkeeping about which index you solve over.

\heading{Step 1 — Count, before you open MATLAB}
\dolabel Write down $n$ (design variables) and $m$ (implicit responses) for both
settings: A is $n=12$, $m=1$; B is $n=2$, $m=7$.\par
\why Direct costs one solve per design variable, the adjoint one per response;
take the smaller \src{slide 22}. The counts swap between the settings, so the
answer swaps --- the method follows the {\itshape shape} of the problem, not the
physics.\par
\rep Setting A: {\bfseries adjoint}. Setting B: {\bfseries direct}. One sentence
each, stating {\itshape both} counts --- the rule is a comparison, not a
preference. Do this before any code; the question asks for it in that order.

\heading{Step 2 — Assemble $A$ and $q$}
\dolabel Build $A$ directly, without forming each $A_e$:
\[
  A_{ii}=\!\!\sum_{\text{pipes at } i}\!\! x_e, \qquad
  A_{ij}=-x_e \ \text{if pipe } e \text{ joins } i,j, \qquad
  A_{ij}=0 \text{ otherwise},
\]
then strike row and column 8, and set $q=e_1$.\par
\why Junction 8 is held at $p_8=0$. Without striking it every row sums to zero,
$A$ is singular and pressure is defined only up to a constant \src{slide 9}.\par
{\color{warn}Trap:} pipes 10 and 11 {\bfseries still count in the diagonals} of
junctions 6 and 7. A grounded neighbour contributes to the diagonal and nothing
else.
\[
\small
A=\begin{bmatrix}
 5.5&-3.0&-2.5& 0 & 0 & 0 & 0\\
-3.0& 6.0&-1.0&-2.0& 0 & 0 & 0\\
-2.5&-1.0& 6.3&-0.8&-2.0& 0 & 0\\
 0 &-2.0&-0.8& 5.8&-1.2&-1.8& 0\\
 0 & 0 &-2.0&-1.2& 4.8& 0 &-1.6\\
 0 & 0 & 0 &-1.8& 0 & 4.1&-0.9\\
 0 & 0 & 0 & 0 &-1.6&-0.9& 3.6
\end{bmatrix}
\]
\chk Rows 1--5 sum to zero. Rows 6 and 7 leave $4.1-1.8-0.9=1.4=x_{10}$ and
$3.6-1.6-0.9=1.1=x_{11}$ --- the conductances to ground. Also $A=A^{\!\top}$. A
row that should sum to zero and does not means a mis-entered pipe. Fix it here,
not three steps later.

\heading{Step 3 — Solve once}
\dolabel $p=A\backslash q$. You should get
\[
  p=(1.09135,\ 0.92901,\ 0.88615,\ 0.70694,\ 0.67990,\ 0.39857,\ 0.40182),
\]
so $f=p_1=1.09135$.\par
\chk Substitute back into any row by hand:
$5.5(1.09135)-3.0(0.92901)-2.5(0.88615)=1.00002\approx q_1$, and row 7 gives
$-0.000001\approx0$. The residuals are your 5-digit rounding, not solver error.\par
\chk Budget: $\sum_e L_ex_e=20.5752\le V=22$. The starting design is feasible ---
say so, it costs one clause.

\heading{Step 4 — Setting A: notice it is self-adjoint {\itshape before} coding}
\dolabel Compare $c$ with $q$. The objective is $f=p_1=e_1^{\!\top}p$, so
$c=e_1$; the flow enters at junction 1, so $q=e_1$. Hence $c=q$, and $A$ is
symmetric.\par
\why Then $A^{\!\top}\lambda=-c$ {\itshape is} the analysis with the load
reversed, so $\lambda=-p$ and no adjoint solve is needed \src{slide 26 --- titled
for this exact setting}. The gradient collapses:
\[
  \frac{df}{dx_e}=\lambda^{\!\top}A_e\,p=-p^{\!\top}A_e\,p=-\left(p_i-p_j\right)^2 .
\]
\rep The chain $c=q$ $\Rightarrow$ $\lambda=-p$ $\Rightarrow$ closed form. This is
the single highest-value sentence in Question 1.

\heading{Step 5 — Get the twelve numbers}
\dolabel $df/dx_e=-(p_i-p_j)^2$ with $p_8=0$. No matrix, no solve --- a
subtraction and a square, twelve times.
\par\smallskip
\small
\begin{tabular}{rlrr|rlrr}
 $e$ & $i$--$j$ & $p_i-p_j$ & $df/dx_e$ & $e$ & $i$--$j$ & $p_i-p_j$ & $df/dx_e$\\[2pt]
\hline
 1&1--2&$0.16234$&$-0.026354$ &  7&4--6&$0.30837$&$-0.095092$\\
 2&1--3&$0.20520$&$-0.042107$ &  8&5--7&$0.27808$&$-0.077328$\\
 3&2--3&$0.04286$&$-0.001837$ &  9&6--7&$-0.00325$&$-0.0000106$\\
 4&2--4&$0.22207$&$-0.049315$ & 10&6--8&$0.39857$&$-0.158858$\\
 5&3--5&$0.20625$&$-0.042539$ & 11&7--8&$0.40182$&$-0.161459$\\
 6&4--5&$0.02704$&$-0.000731$ & 12&3--4&$0.17921$&$-0.032116$\\
\end{tabular}
\normalsize
\par\smallskip
\chk Three routes must agree: the adjoint product $\lambda^{\!\top}A_ep$, this
closed form, and a direct solve. Machine precision ($10^{-16}$) or you have a
bug. A central finite difference agrees only to $\sim10^{-9}$ --- that is its
truncation error and is expected \src{slides 15--16}.\par
{\color{warn}Pipe 9 will look wrong by hand:} $-1.0563\times10^{-5}$ against
$-1.0584\times10^{-5}$, 0.2\% off. The reason is the reason it is small --- $p_6$
and $p_7$ agree to three decimals, so subtracting them destroys your significant
digits. The least accurate entry of a hand check is always the smallest one. Do
not chase it.

\heading{Step 6 — Count the solves you actually performed}
\dolabel Report {\bfseries 2} --- one analysis, one adjoint --- or {\bfseries 1}
if you used $\lambda=-p$. Direct would have cost $1+12=13$.\par
\why One adjoint solve fills a {\itshape row} of the Jacobian, one direct solve a
{\itshape column} \src{slides 18, 21, 22}. Here there is one row and twelve
columns.\par
{\color{warn}Trap:} if you also ran direct as a cross-check, do {\itshape not}
report 13. Report the cost of the method you chose and mention the verification
separately.

\heading{Step 7 — The constraint that needs no analysis}
\dolabel $\partial\!\left(\sum_e L_ex_e\right)\!/\partial x_e=L_e$ --- the pipe
lengths themselves.\par
\why $\sum_e L_ex_e$ is {\itshape explicit} in $x$: no pressure appears, so there
is nothing to differentiate through. The lecture names this very constraint ``the
opposite extreme'' from $c^{\!\top}p$ \src{slide 14}.

\heading{Step 8 — Setting B: recount, then switch method}
\dolabel $n=2$, $m=7$, so {\bfseries direct}: solve
$A\,\partial p/\partial x_e=-A_ep$ for $e=1,2$. Each result is a full column of
the $7\times2$ Jacobian. Total $1+2=3$ solves against the adjoint's $1+7=8$.\par
\why Direct is indexed by the design, and the design is now small --- 14 Jacobian
entries from 3 solves. The lecture works this exact case \src{slide 22}.\par
\chk {\bfseries Row 1 is already known.} Setting A's objective {\itshape is}
$p_1$, so
\[
  \frac{\partial p_1}{\partial x_1}=\frac{df}{dx_1}=-0.026353,
  \qquad
  \frac{\partial p_1}{\partial x_2}=\frac{df}{dx_2}=-0.042106 .
\]
If row 1 does not reproduce entries 1 and 2 of your twelve, one of the two
settings is wrong. Free, and it catches a lot.\par
\chk Signs: pipe 1 feeds the 2--4--6 branch, pipe 2 the 3--5--7 branch, so each
column is positive on its own branch and negative on the other --- and row 1 is
the only row negative in both.

\heading{Step 9 — Setting B's objective}
\dolabel $L_1x_1+L_2x_2$, gradient $(1.118,\,1.118)$. No analysis.\par
\why The mirror of Step 7: in A the objective was implicit and the constraint
free; in B the {\itshape constraints} are implicit and the objective is free. Say
it that way --- the symmetry is the point.

\heading{Step 10 — Read the answer}
\dolabel {\bfseries (1)} Widen {\bfseries pipe 11 (7--8)}, at
$-1.6146\times10^{-1}$, the most negative of the twelve.\par
\why Sign says which way, magnitude says how much \src{slide 25}; and the
magnitude {\itshape is} the squared pressure drop, so the two outlet pipes --- 11
and 10, discharging into the grounded junction --- carry the network's largest
drops ($0.40182$, $0.39857$) and win. Pipe 9 drops $0.00325$ and scores four
orders of magnitude less.\par
\dolabel {\bfseries (2)} Every $df/dx_e\le0$; the largest is
$-1.06\times10^{-5}$.\par
\why {\color{warn}Do not answer ``they all came out negative''.} The question is
why they {\itshape must} be: $A_e=(e_i-e_j)(e_i-e_j)^{\!\top}$ is an outer
product, hence positive semidefinite, hence $(p_i-p_j)^2\ge0$ for every network,
every design, always. Physically, widening a pipe can only make it easier to push
the same flow, so the pump head falls or holds but never rises. That is also why
the volume budget must be active at the optimum.

\heading{Appendix A — what to cite}
\begin{tabular}{@{}p{4.4cm}p{10cm}@{}}
Rao, {\itshape Section 15.4}, p.~705 & Response sensitivity --- the whole of Q1. Assigned before L05 and again before L06, ``now read as the tool that will be used'' \src{slide 29}.\\[3pt]
Rao, {\itshape Section 15.7}, p.~712 & Optimum sensitivity / shadow prices. {\color{warn}Not Q1.}\\[3pt]
L05 slides 1--29 & The derivation chain; slides 5--9 are the small version of this network.\\[3pt]
\texttt{pipe\_sensitivity\_demo.m} & The supplied five-pipe reference script.\\
\end{tabular}

\heading{Appendix B — the distinction most likely to be probed}
\src{slide 27} Two Lagrange multipliers live in this course; Q1 uses only one.
\par\smallskip
\small
\begin{tabular}{@{}p{2.7cm}p{5.0cm}p{5.0cm}@{}}
 & {\bfseries response sensitivity (Q1)} & {\bfseries shadow price (L04, A3 Q2)}\\[3pt]
\hline\\[-6pt]
the question & widen this pipe by one unit --- how far does the pressure move? & move this limit by one unit --- how much better could the best design be?\\[3pt]
the number & $df/dx_e=\lambda^{\!\top}A_ep$ & $df^\star/db=-\lambda$\\[3pt]
you perturb & a design variable & a problem parameter\\[3pt]
you stand & at {\itshape any} design & only at the optimum, after solving\\[3pt]
carried on & the analysis $Ap=q$ & a design limit $g\le0$\\
\end{tabular}
\normalsize
\par\smallskip
Same object, two jobs. Q2's \texttt{lambda.ineqlin} is the right-hand column;
everything in Q1 is the left.

\par\medskip
{\color{warn}Note:} the slides tag results with [Notes \S1.2], [\S2.5], [\S2.6],
[\S5] --- a lecture-notes document {\itshape not} in this repository. Cite slide
numbers, not section numbers you have not seen.

\end{document}
